%% ============================================================
%% Chapter 15: The Compression Fallacy
%% ============================================================

\chapter{The Compression Fallacy}
\label{ch:ch15-compression-fallacy}
\section{An Old Definition, Asked of a New Case}

Part~III closed with a model capsule whose distinctive contribution was to bind a deployed computation's discrete representation to the machinery required to execute it, rather than reporting the representation's size alone. This chapter takes up the general question that binding was built to answer: does a shorter description of something mean that thing has been better understood?

The computer scientist Daniel Hillis once offered a definition of the information content of a pattern: the amount of information in a pattern of bits equals the length of the smallest computer program capable of generating those bits, in substance Kolmogorov complexity restated for a general audience. What this chapter is concerned with is a further claim the definition has often been taken, by others, to support, though it is not one Hillis himself is on record as drawing: that a system which produces a shorter description of something has thereby understood that thing better than a system producing a longer one. The claim is rarely defended directly. It is simply assumed, in the background of arguments about which model of a phenomenon is the better model, which representation of a domain is the more insightful representation, and, increasingly, in discussions of large trained models, which network has learned the most genuine structure rather than having merely memorized surface statistics.

Forth, an unusually terse programming language whose devotees have long pointed to Hillis's definition as evidence of the language's virtue, supplies an unusually clean test case, because Forth's brevity is well documented, uncontroversial, and produced by a mechanism that is easy to state precisely: word-factored Forth code is short because a Forth programmer builds a library of named operations, words, and subsequent code invokes them rather than re-deriving them. A task solved in Forth after such a library exists can be several times shorter, in raw character count, than the same task solved in a general-purpose language without access to an equivalent library. The question is not whether Forth code is short; that is not in dispute. The question is what that shortness is evidence of, and whether it is evidence of the thing Hillis's definition, read the way it is usually read, would have it be evidence of.

This chapter argues that it is not, and that the same gap separating Forth's brevity from Kolmogorov-optimal description separates, more consequentially, model compactness from model understanding wherever the comparison is drawn.

\section{Coding Length and Distinguishing Capacity}

A general framework for measuring representational capacity takes as its basic object not a set of items to be described, nor a description language considered on its own, but the pairing of the two together with a notion of when two items count as the same for the purposes at hand.

\begin{definition}[Ontological triple]
An ontological triple $(X,\sim,T)$ consists of a set $X$, an equivalence relation $\sim$ on $X$ specifying which elements are treated as observationally indistinguishable, and an ontology $T$: a description language, or template hierarchy, determining which descriptions of elements of $X$ are actually reachable, as opposed to merely existing in principle.
\end{definition}

The distinction between what a description language can express in principle and what it can reach in practice, given the templates and operations it has already accumulated, is the hinge the rest of this chapter turns on. Two deficits are defined over this structure, and they measure different things. The coding deficit of an object $x$, written $\delta_T(x)$, is the excess length of the shortest description of $x$ reachable within $T$ over the Kolmogorov-optimal description length of $x$, the length of the absolute shortest program capable of generating $x$, independent of any particular ontology. The distinguishability deficit of $x$, written $\Delta_T(x)$, is a different quantity: let $D_T(x)$ denote the distinguishing capacity of $T$ at $x$, the number of other objects that can be separated from $x$ by descriptions reachable in $T$, and let $D^\ast(x)$ denote the maximum such capacity attainable by any ontology whatsoever. Then $\Delta_T(x) = D^\ast(x) - D_T(x)$. The coding deficit asks how much longer a description is than it needs to be. The distinguishability deficit asks how much distinguishing power has been lost. Nothing in the definitions forces these two quantities to move together, and the relationship between them is the object a theorem is required to establish.

An ontology decomposes as a triple $T = (E_T, M_T, \mathcal{O}_T)$: an encoder $E_T : Q \to \{0,1\}^\ast$ over the semantic quotient $Q = X/\!\sim$, giving each semantic class its reachable description; background machinery $M_T$, the accumulated templates, interpreter, or structure constants required to produce or use those descriptions in the first place; and an operational repertoire $\mathcal{O}_T$, the set of operations computable directly over the encoded representations without first decoding back to $Q$.

\begin{definition}[Three coordinates of a representation]
\label{def:threecoord}
Distinguish three properties: the description magnitude $L_T(q) = |E_T(q)|$; the distinguishing structure $\mathcal{P}_T = Q/\!\sim_T$, where $q \sim_T r \iff E_T(q) = E_T(r)$; and the operational repertoire $\mathcal{O}_T$. These measure different properties of $T$, and equality or improvement in one does not, without an additional coupling assumption supplied from outside the definitions, imply equality or improvement in either of the others:
\[
L_T \not\Rightarrow \mathcal{P}_T, \qquad L_T \not\Rightarrow \mathcal{O}_T, \qquad \mathcal{P}_T \not\Rightarrow \mathcal{O}_T.
\]
\end{definition}

\section{Lossless Compression Does Not Imply Ontological Economy}

Before turning to the machine-learning case, it is worth addressing the objection this apparatus invites immediately: doesn't lossless compression preserve all the information in a representation, so that a shorter lossless encoding must preserve every distinction the longer one did? At the level of static distinguishability this case is straightforward. A ZIP archive can become dramatically shorter than the file it compresses while remaining perfectly invertible: its encoder is injective on $Q$, so $\sim_T$ collapses nothing, and $D_T(q)$ is maximal for every $q$. Lossless compression is accordingly not a counterexample to static distinguishability preservation.

What the ZIP case exposes instead is the third quantity, $\mathcal{O}_T$. That two encoded objects remain distinct, in the static sense $D_T$ measures, does not imply that every relation of interest between their uncompressed contents remains directly computable over the encoded representations themselves. Two compressed files can be told apart trivially, their bitstrings differ, without any geometric or relational property of what they contain being available for comparison until both are decompressed; if decompression is not itself among the operations $\mathcal{O}_T$ makes native, those relations sit outside $\mathcal{O}_T$ even though $D_T$ is already maximal. Description magnitude, distinguishing structure, and operational repertoire are logically independent properties of an ontology, and a lossless archive is the case in which the first two are already exemplary while the third remains, until an inverse transform is paid for, essentially empty.

The converse case rules out the opposite inference. An extremely verbose database, one that stores every field of every record with no compression whatsoever, can distinguish every member of its domain perfectly despite being terribly compressed: $\delta_T(x)$ large, $\Delta_T(x) = 0$. Nobody encountering such a database concludes that its size is evidence against its usefulness, because nobody was using size as a proxy for distinguishing power in the first place; the intuition that a bloated schema is not thereby a poorly understood one is the intuition this chapter is asking to be extended to the cases, sparse latent codes, monosemantic feature dictionaries, where the opposite inference is currently made routinely.

\section{Compression Relative to What?}

A further assumption is buried in any claim that a representation is short: short relative to what reference machinery. Kolmogorov complexity has its invariance theorem, which guarantees that the choice of universal machine affects $L^\ast(x)$ by at most an additive constant independent of $x$, but that guarantee is asymptotic, and practical comparisons between representational systems are made nowhere near the regime in which an additive constant can be safely ignored. Writing $L_{\mathrm{visible}}(q) = L_T(q)$ for the quantity ordinarily reported and $L_{\mathrm{system}}(q) = L(M_T) + L_T(q)$ for the full cost,
\[
L_{T_2}(q) < L_{T_1}(q) \quad \not\Rightarrow \quad L(M_{T_2}) + L_{T_2}(q) < L(M_{T_1}) + L_{T_1}(q).
\]
Forth makes the point vivid: the token \texttt{FOO} is three characters, but the definition of \texttt{FOO} together with the words it depends on can represent thousands of characters of accumulated machinery. Reporting the visible length without the machinery is not measuring a smaller quantity than the full system cost; it is measuring a different, incomplete one. The same separation generalizes past Forth without alteration: a short latent code in a trained autoencoder can depend on a decoder with billions of parameters, and reporting the latent code's length while omitting the decoder's is the identical move under a different name.

\begin{proposition}[Machinery Displacement]
\label{prop:machinery-displacement}
Let $Q = X/\!\sim$ be finite. There exists an ontology $T_2 = (E_2,M_2)$ in which every $q \in Q$ has a fixed-width description of length $\lceil \log_2 |Q| \rceil$, while the machinery required to interpret those descriptions contains the corresponding lookup structure. Consequently, reducing the per-object description length does not by itself establish that representational complexity has been eliminated rather than displaced into background machinery.
\end{proposition}

Enumerate $Q = \{q_0,\ldots,q_{n-1}\}$ and define $E_2(q_i) = \mathrm{bin}_{\lceil \log_2 n \rceil}(i)$, with $M_2$ containing the correspondence between indices and the classes they denote. Every encoded object then has the same short description length, while interpreting those descriptions presupposes the correspondence stored in $M_2$. The question worth asking of any short description is accordingly not how short is it, but what machinery must already exist for it to recover the object it names. A shell command invokes an entire operating system. A Forth word invokes its dictionary. A latent vector invokes a decoder. A theorem invokes its definitions and lemmas.

\section{Description-Length Non-Identification}

\begin{proposition}[Description-Length Non-Identification]
\label{prop:nonident}
Let $Q = X/\!\sim$ be finite with $|Q| \geq 3$, and let $q \in Q$. Then: (i) for every pair $(\ell,k)$ with $\ell \geq 1$ and $0 \leq k \leq |Q|-1$, there exists an ontology $T$ with $L_T(q) = \ell$ and $D_T(q) = k$; and (ii) for each relation $R \in \{<,=,>\}$, there exist $T_1, T_2$ with $L_{T_1}(q) < L_{T_2}(q)$ and $D_{T_1}(q)\, R\, D_{T_2}(q)$.
\end{proposition}

Let $n = |Q|$. For (i), choose a subset $C \subseteq Q$ containing $q$ with $|C| = n-k$. Define the encoder so every element of $C$ receives the same description $s$ of length $\ell$, and assign every element of $Q \setminus C$ a description distinct from $s$ and from every other such description, which is always possible since $\{0,1\}^\ast$ is infinite. The encoder-induced class of $q$ is then exactly $C$, so $D_T(q) = k$ and $|E_T(q)| = \ell$. Statement (ii) follows by applying (i) twice with $\ell_1 < \ell_2$ and $k_1, k_2$ chosen so that $k_1\, R\, k_2$.

A fully worked instance, fixing $Q = \{a,b,c\}$ and comparing five encoders at $q=a$, makes the construction checkable directly. A total-collapse encoder gives $(L,D) = (3,0)$; separating $a$ alone from $b,c$ merged gives $(1,2)$; grouping $a$ with $b$ gives $(2,1)$; an injective two-character encoding gives $(2,2)$; and a padded injective four-character encoding gives $(4,2)$. Comparing the third and fourth: equal length, unequal capacity. Comparing the second and first: shorter length, strictly higher capacity. Comparing the third and fifth: shorter length, lower capacity, the case ordinary usage means by ``compression loses information,'' present here as one populated cell among several rather than the only possibility. Comparing the fourth and fifth: shorter injective encoding against a longer injective one, identical capacity. Every cell of the ordering table is realized by an explicit construction at the same fixed point; none is left as an existence claim.

\[
\boxed{\text{compression magnitude does not determine compression topology}}
\]

\begin{proposition}[Dynamic non-identification]
\label{prop:dynamic}
Let $(T_t)_{t=0}^m$ be a sequence of ontologies over a fixed finite $Q$, and $q \in Q$. Write $\Delta L_t = L_{T_{t+1}}(q) - L_{T_t}(q)$ and $\Delta D_t = D_{T_{t+1}}(q) - D_{T_t}(q)$. Without additional constraints linking description length to encoder-induced partitions, $\operatorname{sgn}(\Delta L_t) \not\Rightarrow \operatorname{sgn}(\Delta D_t)$.
\end{proposition}

Apply Proposition~\ref{prop:nonident} independently at times $t$ and $t+1$, choosing $\ell_{t+1} < \ell_t$ while choosing $k_t, k_{t+1}$ to realize any desired ordering between them. Differentiating an unidentified relationship with respect to time does not identify it. A compression-progress account that reads $\Delta L_t < 0$ as evidence of improving understanding, or a temporary $\Delta L_t > 0$ as evidence of regress, commits Proposition~\ref{prop:nonident}'s error once per time step, not avoided by moving to a derivative.

\section{The Proxy Theorem as a Conditional Result}

A previously proposed result, the Deficit Proxy Theorem, is worth examining as a candidate additional-assumption case of Proposition~\ref{prop:nonident}, the kind of result that pins down one point on an otherwise unidentified surface by adding a specific, checkable assumption. Under an assumption of faithful encoding, that every distinction a description language can express costs at least one bit, the theorem states that the distinguishability deficit is bounded above by a constant multiple of the coding deficit, with the special case that a zero coding deficit implies a zero distinguishability deficit: an ontology achieving the Kolmogorov-optimal description of an object would thereby also achieve the maximum possible distinguishing capacity for it.

Closer inspection shows even this special case is not secured by the argument given for it. The step from Kolmogorov-optimality to the claimed conclusion treats generating an object cheaply and distinguishing that object from an ambient domain as though they were the same achievement; they are not, and closing the gap would require an additional condition tying the optimal description length to a maximally distinguishing encoding specifically, a condition the source result does not supply. The theorem is presented here as a previously proposed conditional result under examination rather than as an established point this chapter's own argument depends on. The converse, that a distinguishability deficit of zero forces a coding deficit of zero, fares no better; it is stated only as a conjecture in its source, resting on the further, unargued assumption that optimal coding uses every distinction available to it.

The asymmetry matters because the popular reading of Hillis's definition uses the coding deficit as a graded, continuous stand-in for the distinguishability deficit across its whole range, treating even the strongest available version of the Deficit Proxy Theorem as though it licensed that continuous use. It would not license it even if its own special case were secure: an inference from perfect compression to perfect distinguishing capacity is not an inference that a further reduction in an already-imperfect coding deficit constitutes evidence of a further reduction in distinguishing capacity.

\section{Forth and Reachable Complexity}

Forth supplies reachable complexity's first concrete instance: the shortest description of an object expressible from a description language's current template hierarchy, which may be strictly longer than the Kolmogorov-optimal description. A Forth programmer accumulates, over the life of a project, a library of named words, short, composable operations built from more primitive ones, each capturing an operational distinction the programmer has found useful to keep separately addressable. Code written after such a library exists is short because it is reachable, cheaply, from that library, not because it approaches the description length that would be achieved by an ontology with no library at all, starting from nothing.

This inverts the naive reading of Hillis's definition as applied to Forth. The naive reading takes Forth's brevity as evidence that Forth programs approach the information-theoretic floor of the tasks they solve. The reachable-complexity reading takes the same brevity as evidence of something else: that a well-factored library of admissible operations was available to be reused. Two Forth programmers solving the same task with different accumulated word libraries will produce differently short code for it, and the difference in length reports which library was available, not which programmer better understood the task. A companion result states the point in general form: no valid compression mapping can be constructed without prior identification of the invariant structure of the space being compressed, and that identification is itself achieved only through a prior process of expansion, exploration, iteration, and discovery of which structure is invariant in the first place, developed at length in Chapter~\ref{ch:ch16-compression-after-expansion}. Forth's terseness is compression in exactly this downstream sense: it presupposes the expansion that built the word library, and reports that expansion's prior existence rather than the Kolmogorov-optimal floor of the task the short code happens to solve.

None of this counts against Forth. A large, well-factored word library achieving low coding deficit relative to itself, while the Kolmogorov-optimal length remains a separate and largely inaccessible quantity, is exactly what a mature, high-distinguishing-capacity representational system should look like. What the coding deficit cannot do, without the further, unproved direction of the Proxy Theorem, is stand in for a claim about how much of the task's true distinguishing structure the library actually preserves.

\section{The Same Move in Machine-Learned Representations}

The informal assumption this chapter has been arguing against is closer to background theory in contemporary discussions of trained models, where a shorter, sparser, or more compressed internal representation is routinely treated as evidence of deeper or more genuine understanding. The assumption's most explicit and most influential statement extends compression progress well past a technical modeling virtue and treats it as the underlying currency of understanding itself, a unifying account of curiosity, creativity, and aesthetic interest in which subjective beauty and interestingness are identified with the rate at which an observer's compression of its data is currently improving. If the relationship between coding deficit and distinguishability deficit is only proven in one direction and only at the extreme, then a theory that identifies the felt sense of understanding, or beauty, or interest with compression progress at every point along the curve, rather than only at the point where compression is complete, is claiming a great deal more than the formal relationship between the two quantities actually licenses.

Two examples make the pattern concrete. The beta-VAE architecture increases a single compression pressure, a weighting term on the model's latent-code cost, specifically in order to produce latent representations described as more interpretable and as more faithfully capturing a domain's independent underlying factors of variation, with the increase in compression pressure offered as the mechanism by which the improvement is obtained. Sparse-autoencoder interpretability work on language models motivates the move to sparse, low-dimensional dictionaries of features in comparable terms: individual neurons are described as polysemantic and therefore poorly understood, while the sparse, more compressed feature directions recovered by the autoencoder are described as monosemantic and correspondingly more interpretable. In neither case is the inference from compression to understanding stated as an unproven assumption requiring a caveat; in both, it is the stated motivation for the method.

The apparatus above applies without modification, and the case exposes a further difficulty the Forth case does not. The Deficit Proxy Theorem's special case depends on faithful encoding, that every distinction costs at least one unit of representational capacity. Distributed and superposed representations, of the kind large trained networks are now widely understood to use, in which a single direction in activation space participates in encoding more than one feature, and a single feature is encoded across more than one direction, are precisely representations for which faithful encoding need not hold. Where it fails, even the weakened special case no longer applies, and a reduction in a representation's measured description length carries no guaranteed relationship to its distinguishing capacity at all, in either direction. The inference from compression to understanding is, in this setting, not merely unsupported past some extreme; it is not obviously licensed anywhere along the range where trained models actually operate.

\section{Overcompression and Catastrophic Identification}

It is worth naming directly what happens at the point where the divergence between the two deficits becomes irreversible. Suppose $x_1 \neq x_2$ but an encoding maps both to the same representation. Compression has crossed an epistemic boundary at that pair: unless information external to the encoding remains available, the distinction cannot subsequently be reconstructed. Every compression scheme implicitly asserts that the distinctions collapsed within its induced equivalence classes do not matter, whether or not that assertion was made deliberately. The question worth asking of any compressed representation is accordingly not how much was compressed but which differences the compression declared equivalent, a question about $\mathcal{P}_T$, not about $L_T$.

\begin{definition}[Task-relative equivalence and overcompression]
For a task $\tau : X \to Y$, define $x \sim_\tau y \iff \tau(x) = \tau(y)$. An encoding $E$ preserves the distinctions $\tau$ requires when $E(x) = E(y) \implies x \sim_\tau y$, equivalently $\sim_E\, \subseteq\, \sim_\tau$. $E$ overcompresses relative to $\tau$ if there exist $x,y$ with $E(x) = E(y)$ but $\tau(x) \neq \tau(y)$.
\end{definition}

A statistic $S(X)$ can throw away enormous quantities of raw information while retaining everything necessary for inference about a specified parameter $\theta$; if $S$ is sufficient for $\theta$, it does not overcompress relative to $\theta$, however much it overcompresses relative to $X$ itself. Compression, on this reading, is not opposed to understanding; it is meaningful only relative to a prior decision about which distinctions must survive. The failure this section names is not compression itself but overcompression relative to some further task the representation is silently expected to also serve: a representation with $\sim_E \subseteq \sim_\theta$ mistaken for one with $\sim_E \subseteq \sim_X$.

\section{Curricula Must Widen}

A further difficulty for any theory that identifies understanding, or interest, or progress with monotonic compression is that the environments in which learning actually occurs are not found in a state of nature; they are built, and built deliberately simple, for reasons that have nothing to do with the learner's understanding having reached some genuine informational floor. Children are not raised inside the full chaos of an unfiltered natural environment; they are raised inside environments deliberately constrained to present regularities slowly, precisely because a learner who cannot yet compress much of anything will fail to learn from an environment offering nothing but unstructured, simultaneous complexity. The same engineering shows up outside pedagogy for reasons that are economic rather than developmental: standardized building materials constrain the built environment to a narrow vocabulary of regular shapes because mass production rewards uniformity, not because simple, rectilinear spaces are closer to the true shape of anything; roads are built straight because the vehicles traveling them have limited suspension travel, not because the terrain they cross is actually simple.

In each case, the low entropy of the environment is a scaffold fitted to the current, limited distinguishing capacity of whatever is operating within it, not a report on the true complexity of the domain. A scaffold is only useful for as long as the learner has not yet outgrown it. Once a learner is no longer acquiring anything new from a fixed training environment, further progress requires increasing the entropy of what is presented, widening $D_T(x)$ toward the distinctions the scaffold had been withholding, not continuing to compress what is already compressed as far as it usefully can be. A curriculum that only ever simplifies, and never widens, terminates in a plateau rather than in understanding. This is difficult to reconcile with a theory that identifies the drive toward understanding with compression progress at every point along the learner's trajectory, since the trajectory that theory describes is one of steadily increasing compression, while the trajectory that actually produces continued learning alternates compression with deliberately reintroduced complexity.

A related and sharper difficulty appears within problem-solving itself, on timescales far shorter than a curriculum. Solving a problem often requires making it more complicated for a while before it can be made simple: introducing an auxiliary variable that has no simplifying role on its own but opens a path to one, or embedding a problem in a higher-dimensional space where a solution becomes visible before projecting the result back down. A compression-progress account, tracking description length as it evolves along the solver's trajectory, would register this intermediate expansion as a local reversal precisely at the point where the solver is doing the genuine work that will shortly produce the solution.

Developmental research on the explore-exploit trade-off supplies independent, empirically grounded support for the same structural claim: the distinctive length of human childhood is itself argued to be an evolved solution to a tension between exploring broadly and exploiting what is already known, with early cognition characterized by wide, high-entropy hypothesis search progressively narrowed only as the demands of goal-directed exploitation increase. On this account childhood is not an inefficient, not-yet-compressed version of adult cognition working its way toward an adult's compression; it is a distinct cognitive strategy favoring the preservation of distinctions a purely exploitative learner would discard, for exactly as long as the developmental trade-off favors exploration over efficiency.

\section{Understanding Under Intervention}

Measuring distinguishing capacity directly rather than inferring it from description length is still, as stated so far, a claim about a static quantity computed once against a fixed domain. A stronger and more operational version asks not how many distinctions a representation preserves at rest, but which distinctions survive when the domain is deliberately perturbed. Write $D_T(x \mid I)$ for the distinguishing capacity of $T$ at $x$ after an intervention $I$, and define the intervention robustness $R_T(x,I) = D_T(x\mid I) - D_T(x)$. Two representations with identical description length, or even identical distinguishing capacity measured before any intervention, can behave completely differently under $I$:
\[
L_{T_1}(x) = L_{T_2}(x) \ \text{and}\ D_{T_1}(x) = D_{T_2}(x) \ \not\Rightarrow\ D_{T_1}(x\mid I) = D_{T_2}(x\mid I).
\]
A geometric construction examined in Chapter~\ref{ch:ch18-clifford-fhe}, in which rotational equivariance follows from an algebraic identity of the geometric product rather than being an empirical regularity discovered by training on many rotated examples, supplies a motivating case that is unusually clean because the claim being tested is architectural rather than statistical: intervention robustness under rotation is there provable from the representation's algebra, not merely measured empirically. A representation's distinguishing capacity is revealed by which distinctions survive admissible transformations of the domain, not merely by how many bits are required to encode the domain's resting state.

\section{What Should Be Measured Instead}

Proposition~\ref{prop:nonident} establishes that coding length and distinguishing capacity can come apart across their entire joint range except at one identified point, and the Forth and machine-learning cases show the gap is not merely possible but ordinary. A practice that measures the first while reporting conclusions about the second is measuring the wrong quantity almost everywhere it is applied. None of this is an argument that compression is irrelevant to understanding. A system with a smaller coding deficit is, all else equal, a more efficient one, and efficiency is its own legitimate goal. The correction is a separation of that virtue from claims about understanding, which should be evaluated against the distinguishability deficit directly wherever the two quantities can be independently assessed: does the representation actually preserve the ability to separate the objects, cases, or outcomes that matter for the task, rather than merely occupying fewer bits while doing so.

A useful example already appears in evaluation practice for text summarization. Systems designed to produce very short summaries confront the obvious fact that brevity can preserve a coarse overview while eliminating distinctions present in the source; some summarization work supplements automatic overlap measures with separate human judgments of completeness and cohesion, rather than treating the automatic score as sufficient. The same system's central mechanism supplies a more literal instance of this chapter's apparatus: a summarizer that selects short introductory sentences and uses them as pointers into the rest of the source document exhibits, in the reach a pointer sentence makes available under a trained selection network, exactly the non-identification Proposition~\ref{prop:nonident} describes: a short pointer sentence can have a large reach given a sufficiently capable selector, and a longer one can have almost none. The visible sentence is $L_T$; the trained selector is $M_T$; what the system is actually engineered to get right is the reach, which the sentence's own length does not report. Three routes past the difficulty recur throughout this chapter: ask what library of prior, admissible distinctions a representation's shortness is reachable from, and whether that library actually covers the distinctions the task requires; ask whether a representation's apparent compactness survives an audit for superposition and faithful encoding before treating that compactness as evidence of anything beyond itself; and, most demanding of the three, perturb the domain and ask which distinctions survive the perturbation, rather than trusting a static count taken at rest.
